\begin{statement}{A.29}{Proved}{At most one odd-prime stable genus-three packet remains}
Let $\ell$ be odd and let
$C:v^2=w^7-w$ with its two cusps removed. Use an induced
original packet action and a permitted open subgroup
$\Delta_0$ as in A.24 and A.28. If its ordinary first
Johnson homomorphism is nonzero, then
\[
                    P_{C,\ell}^{\,s_b(\Delta_0)}=I_{C,\ell}.
                                                        \tag{A.29.1}
\]
No invertibility hypothesis on any first-Johnson matrix
is needed.

For each odd $\ell$, there is at most one original
pro-$\ell$ packet class, up to conjugation by original
cusp inertia, for which the induced first Johnson
homomorphism vanishes after finite descent. Every other
original packet class satisfies (A.29.1) for every
permitted open subgroup. In particular, along every
nontrivial Kummer direction in A.24, at most one
parameter can remain undecided on $C$.

This does not assert that the exceptional class exists,
identify its tangent parameter, or determine its full
fixed subgroup. Vanishing here means vanishing on an
open subgroup of $\Delta$, rather than just on the
possibly smaller group $\overline{[G_L,G_L]}$.
\end{statement}
\begin{proof}
\paragraph{A nonzero matrix is enough after the monodromy reduction.}
Write the ordinary graded Lie algebra as
$\mathbb L(kt\oplus V)$, where $k=\mathbb Q_\ell$,
$t$ is the selected cusp class, and $V$ is the
six-dimensional negative hyperelliptic eigenspace.
For every $\delta\in\Delta_0$, A.24 gives
\[
 D_\delta(t)=0,\qquad
 D_\delta(v)=[t,A_\delta v],\qquad
                  A_\delta J=JA_\delta^T,
                                                        \tag{A.29.2}
\]
where $\Omega=\sum J_{ij}v_iv_j$ in the tensor algebra
and $J$ is nonsingular and skew. Define the quadratic
Lie element
\[
                     \Omega_{A_\delta}
                              =\sum(A_\delta J)_{ij}v_iv_j.
\]
It is a Lie element because $A_\delta J$ is skew, and
it is nonzero precisely when $A_\delta\ne0$.
The same tensor expansion that gives the peripheral
constraint in A.24 gives the exact identity
\[
                   D_\delta(\Omega)
                             =[t,\Omega_{A_\delta}].
                                                        \tag{A.29.3}
\]
The terms with a $t$ between two $V$ letters cancel
by (A.29.2).

We claim that, whenever $A_\delta\ne0$,
\[
 \ker D_\delta\ \cap\ \mathbb L_k(t,\Omega)=kt.
                                                        \tag{A.29.4}
\]
If $\Omega_{A_\delta}=a\Omega$ for $a\ne0$, the
restriction is $a\operatorname{ad}(t)$, so this is
the free-generator centralizer calculation. Otherwise
the two quadratic tensors $\Omega,\Omega_{A_\delta}$
are linearly independent. Their tensor algebra
embeds into $T(V)$ by grouping factors in pairs,
and adjoining the independent letter $t$ still
gives a free associative algebra. There is therefore
a homomorphism on this algebra that fixes $t,\Omega$
and sends $\Omega_{A_\delta}$ to $\Omega$.
For every $u\in\mathbb L(t,\Omega)$ it sends
$D_\delta(u)$ to $[t,u]$, by (A.29.3) and the
derivation rule. Thus $D_\delta(u)=0$ implies
$[t,u]=0$, proving (A.29.4).

Let $u$ be an actual fixed pro-$\ell$ element.
If the first local monodromy alternative in A.28
occurs, that statement already proves (A.29.1).
Otherwise its two actual second-degree derivations
put the first nonzero class of $u$ in
$\mathbb L(t,\Omega)$. These derivations come from
the normal commutator group used there, contained
in $\Delta_0$. We may also apply the first
derivation of any $\delta\in\Delta_0$ with
$A_\delta\ne0$. By (A.29.4), the leading class
is a multiple of $t$. Its coefficient lies in
$\mathbb Z_\ell$ because $t$ is primitive in
integral homology. Subtract the corresponding
actual fixed cusp power. A nontrivial remainder
would have a leading class of degree at least
two in the same common kernel $kt$, which is
impossible. This proves (A.29.1).

\paragraph{The exceptional original packet class is unique if it exists.}
The first Johnson map is an additive continuous
homomorphism on any permitted $\Delta_0$, since
this group acts trivially on integral homology.
Its target is a torsion-free finite-dimensional
$\mathbb Q_\ell$-space. If it vanishes on an open
subgroup, it vanishes on the whole group:
each element has a positive power in that
subgroup. Hence the zero-Johnson condition is
independent of further finite descent.

Suppose two original packet pro-$\ell$ classes
have this condition. Choose a common finite
descent field and a common open subgroup
$\Delta_0$ for the two induced actions.
Let $\kappa:G_K\to\mathbb Z_\ell(1)$ be their
original cusp Kummer difference. The actual
cusp ramification index for the stable $S_3$
cover is three. Exactly as in A.24, the two
first-Johnson matrices on the common subgroup
differ by
\[
                        \frac{\kappa(\delta)}3\,1_V.
                                                        \tag{A.29.5}
\]
At $\ell=3$ the common finite descent ensures
that the numerator is divisible by three in
the integral cusp action. At other primes
three is a unit. In all cases zero on both
sides forces $\kappa|_{\Delta_0}=0$.

The restriction of $\kappa$ to $\Delta$ is an
additive character into a torsion-free group.
Vanishing on its open subgroup $\Delta_0$
therefore implies $\kappa|_\Delta=0$.
It descends to a cyclotomic cocycle on
$G_K/\Delta$. The odd-prime calculation in
A.6 gives
\[
              H^1_{\mathrm{cts}}
                 (G_K/\Delta,\mathbb Z_\ell(1))=0.
                                                        \tag{A.29.6}
\]
Thus $\kappa$ is a cusp coboundary and the
two original pro-$\ell$ packet classes coincide.
This proves the asserted uniqueness, without
assuming any value for the actual matrices.

Finally a nontrivial Kummer direction used
in A.24 has nonzero restriction to every
permitted open subgroup of $\Delta$, by
the local Frobenius--tame commutator there.
Two distinct parameters on that ray cannot
both have zero first Johnson map, by
(A.29.5). This gives the one-parameter
conclusion as well.
\end{proof}
